\documentclass[noauthor,noinstructornotes]{ximera}

\title{Bookselling}
\author{Kenneth Berglund}

\begin{document}
\begin{abstract}
    Answering questions about a bookselling business.
\end{abstract}
\maketitle

Let \(C\) be the function that gives the average cost per book 
\(C\left(x\right)\), in dollars, when using an online store to print
\(x\) copies of a self-published paperback book. Here is a graph of 
\(c\left(x\right)=\frac{120+4x}{x}\). 

\begin{center}
    \includegraphics[width=20cm]{bookselling.png}
\end{center}

\begin{enumerate}
    \item What is the approximate cost per book when 50 books are 
    printed? 100 books?
    \item The author plans to charge \$8 per book. About how many
    should be printed to make a profit?
    \item What is the value of \(C(x)\) when \(x = \frac{1}{2}\)? how
    does this relate to the context?
    \item What does the end behavior of the function \(C\) say about
    the context. Write down your answer, then share with a partner.
\end{enumerate}

\begin{instructorNotes}
For the third part, invite the discussion of the ``half-book''. Does
selling a half-book make sense? What about a book that's half as long?
\end{instructorNotes}
\end{document}
